% Part: first-order-logic
% Chapter: syntax-and-semantics
% Section: assignments

\documentclass[../../../include/open-logic-section]{subfiles}

\begin{document}

\olfileid{fol}{syn}{ass}

\olsection{变元指派}

\begin{explain}
变元指派~$s$ 为\emph{每个}变元指定一个值——而变元有无穷多个。
当然，这并非逻辑上必需。
我们要求变元指派为所有变元赋值，仅仅因为这样会使事情简单得多。
项~$t$ 的值，以及公式~$!A$ 是否在某个结构中相对于~$s$ 得到满足，
仅仅依赖于指派~$s$ 对~$t$ 中的变元以及~$!A$ 的自由变元所作的赋值。
这正是下面两个命题的内容。
为了使“依赖于”这一想法变得精确，我们将证明：
任何两个在~$t$ 的所有变元上一致的变元指派，都给出相同的值；
并且，如果两个变元指派在~$!A$ 的所有自由变元上一致，
那么~$!A$ 相对于其中一个得到满足，当且仅当它相对于另一个得到满足。
\end{explain}

\begin{prop}\ollabel{prop:valindep}
如果项~$t$ 中的变元都在 $x_1$, \dots,~$x_n$ 之中，并且
$s_1(x_i) = s_2(x_i)$ 对 $i = 1$, \dots,~$n$ 成立，那么 $\Value{t}{M}[s_1]
= \Value{t}{M}[s_2]$。
\end{prop}

\begin{proof}
对~$t$ 的复杂度进行归纳。对于基始情况，$t$ 可以是常项符号
或变元 $x_1$, \dots,~$x_n$ 之一。
若 $t = c$，则 $\Value{t}{M}[s_1] = \Assign{c}{M} = \Value{t}{M}[s_2]$。
若 $t = x_i$，由本命题假设有 $s_1(x_i) = s_2(x_i)$，
因此 $\Value{t}{M}[s_1] = s_1(x_i) = s_2(x_i) = \Value{t}{M}[s_2]$。

对于归纳步骤，假设 $t = \Atom{f}{t_1, \dots, t_k}$
且该断言对 $t_1$, \dots, $t_k$ 成立。则
\begin{align*}
  \Value{t}{M}[s_1] & = \Value{\Atom{f}{t_1, \dots, t_k}}{M}[s_1] \\
  & = \Assign{f}{M}(\Value{t_1}{M}[s_1], \dots, \Value{t_k}{M}[s_1]).
\intertext{对 $j = 1$, \dots,~$k$，$t_j$ 中的变元都在 $x_1$, \dots,~$x_n$ 之中。
根据归纳假设，$\Value{t_j}{M}[s_1] = \Value{t_j}{M}[s_2]$。于是，}
\Value{t}{M}[s_1] & = \Value{\Atom{f}{t_1, \dots, t_k}}{M}[s_1] \\
  & = \Assign{f}{M}(\Value{t_1}{M}[s_1], \dots, \Value{t_k}{M}[s_1]) \\
  & = \Assign{f}{M}(\Value{t_1}{M}[s_2], \dots, \Value{t_k}{M}[s_2]) \\
  & = \Value{\Atom{f}{t_1, \dots, t_k}}{M}[s_2] = \Value{t}{M}[s_2].
\end{align*}
\end{proof}

\begin{prop}\ollabel{prop:satindep}
如果 $!A$ 中的自由变元都在 $x_1$, \dots,~$x_n$ 之中，并且
$s_1(x_i) = s_2(x_i)$ 对 $i = 1$, \dots,~$n$ 成立，那么 $\Sat{M}{!A}[s_1]$
当且仅当 $\Sat{M}{!A}[s_2]$。
\end{prop}

\begin{proof}
我们对 $!A$ 的复杂度进行归纳。对于基始情况，即 $!A$ 为原子公式时，$!A$ 可以是：
\iftag{prvTrue}{$\ltrue$，}{}
\iftag{prvFalse}{$\lfalse$，}{}
$\Atom{R}{t_1, \dots, t_k}$（其中 $R$ 是一个 $k$ 元谓词符号，$t_1$, \dots,~$t_k$ 是项），或者 $\eq[t_1][t_2]$（其中 $t_1$ 和~$t_2$ 是项）。
对于后两种情况，我们只证明双条件句的前向方向，因为反向的证明是对称的。

\begin{enumerate}
\tagitem{prvTrue}{%
  \indcase{!A}{\ltrue}{同时有 $\Sat{M}{\indfrm}[s_1]$ 和
    $\Sat{M}{\indfrm}[s_2]$。}}{}

\tagitem{prvFalse}{%
  \indcase{!A}{\lfalse}{同时有 $\Sat/{M}{!A}[s_1]$ 和
    $\Sat/{M}{!A}[s_2]$。}}{}

\item
  \indcase{!A}{\Atom{R}{t_1, \ldots, t_k}}{设
    $\Sat{M}{\indfrm}[s_1]$。则
    \[
    \langle \Value{t_1}{M}[s_1], \ldots, \Value{t_k}{M}[s_1] \rangle
    \in \Assign{R}{M}.
    \]
    对 $i = 1$, \dots,~$k$，由 \olref{prop:valindep} 有 $\Value{t_i}{M}[s_1] =
    \Value{t_i}{M}[s_2]$。因此我们也有
    $\langle \Value{t_i}{M}[s_2], \ldots, \Value{t_k}{M}[s_2] \rangle
    \in \Assign{R}{M}$，从而 $\Sat{M}{\indfrm}[s_2]$。}
\item
  \indcase{!A}{\eq[t_1][t_2]}{假设 $\Sat{M}{\indfrm}[s_1]$。
    则 $\Value{t_1}{M}[s_1] = \Value{t_2}{M}[s_1]$。于是，
    \begin{align*}
      \Value{t_1}{M}[s_2] & = \Value{t_1}{M}[s_1]
      & \text{（由 \olref{prop:valindep}）} \\
      & = \Value{t_2}{M}[s_1]
      & \text{（因为 $\Sat{M}{\eq[t_1][t_2]}[s_1]$）}\\
      &= \Value{t_2}{M}[s_2]
      & \text{（由 \olref{prop:valindep}），}
    \end{align*}
    所以 $\Sat{M}{\eq[t_1][t_2]}[s_2]$。}
\end{enumerate}

现在假设对所有复杂度低于~$!A$ 的公式 $!B$，都有 $\Sat{M}{!B}[s_1]$ 当且仅当 $\Sat{M}{!B}[s_2]$。归纳步骤根据~$!A$ 的主联结词分情况讨论。
在每种情况下，我们只证明双条件句的前向方向；反向的证明是对称的。除量词的情况外，我们将归纳假设应用于~$!A$ 的子公式~$!B$。~$!B$ 的自由变元包含在~$!A$ 的自由变元之中。因此，如果 $s_1$ 和 $s_2$ 在~$!A$ 的自由变元上一致，那么它们在~$!B$ 的自由变元上也一致，从而归纳假设适用于~$!B$。

\begin{enumerate}
\tagitem{defNot}{}{%
  \iftag{probNot}{%
    \indcase!{!A}{\lnot !B}{}}{%
    \indcase{!A}{\lnot !B}{若 $\Sat{M}{\indfrm}[s_1]$，则
      $\Sat/{M}{!B}[s_1]$，故由归纳假设，
      $\Sat/{M}{!B}[s_2]$，因此 $\Sat{M}{\indfrm}[s_2]$。}}}

\tagitem{defAnd}{}{%
  \iftag{probAnd}{%
    \indcase!{!A}{!B \land !C}{}}{%
    \indcase{!A}{!B \land !C}{若 $\Sat{M}{\indfrm}[s_1]$，则
      $\Sat{M}{!B}[s_1]$ 且 $\Sat{M}{!C}[s_1]$，故由归纳假设，
      $\Sat{M}{!B}[s_2]$ 且 $\Sat{M}{!C}[s_2]$。因此，
      $\Sat{M}{\indfrm}[s_2]$。}}}

\tagitem{defOr}{}{%
  \iftag{probOr}{%
    \indcase!{!A}{!B \lor !C}{}}{%
    \indcase{!A}{!B \lor !C}{若 $\Sat{M}{\indfrm}[s_1]$，则
      $\Sat{M}{!B}[s_1]$ 或 $\Sat{M}{!C}[s_1]$。由归纳假设，
      $\Sat{M}{!B}[s_2]$ 或 $\Sat{M}{!C}[s_2]$，故 $\Sat{M}{\indfrm}[s_2]$。}}}

\tagitem{defIf}{}{%
  \iftag{probIf}{%
    \indcase!{!A}{!B \lif !C}{}}{%
    \indcase{!A}{!B \lif !C}{若 $\Sat{M}{\indfrm}[s_1]$，则
    $\Sat/{M}{!B}[s_1]$ 或 $\Sat{M}{!C}[s_1]$。由归纳假设，
    $\Sat/{M}{!B}[s_2]$ 或 $\Sat{M}{!C}[s_2]$，故 $\Sat{M}{\indfrm}[s_2]$。}}}

\tagitem{defIff}{}{%
  \iftag{probIff}{%
    \indcase!{!A}{!B \liff !C}{}}{%
    \indcase{!A}{!B \liff !C}{若 $\Sat{M}{\indfrm}[s_1]$，则或者
      $\Sat{M}{!B}[s_1]$ 且 $\Sat{M}{!C}[s_1]$，或者 $\Sat/{M}{!B}[s_1]$ 且
      $\Sat/{M}{!C}[s_1]$。由归纳假设，或者
      $\Sat{M}{!B}[s_2]$ 且 $\Sat{M}{!C}[s_2]$，或者 $\Sat/{M}{!B}[s_2]$
      且 $\Sat/{M}{!C}[s_2]$。在任何一种情况下，都有 $\Sat{M}{\indfrm}[s_2]$。}}}

\tagitem{defEx}{}{%
  \iftag{probEx}{%
    \indcase!{!A}{\lexists[x][!B]}{}}{%
    \indcase{!A}{\lexists[x][!B]}{若 $\Sat{M}{\indfrm}[s_1]$，则存在
      某个 $m \in \Domain{M}$ 使得 $\Sat{M}{!B}[\Subst{s_1}{m}{x}]$。
      设 $s_1' = \Subst{s_1}{m}{x}$ 和 $s_2' =
      \Subst{s_2}{m}{x}$。~$!B$ 的自由变元在 $x_1$,
      \dots, $x_n$ 和 $x$ 之中。$s_1'(x_i) = s_2'(x_i)$，因为 $s_1'$ 和
      $s_2'$ 分别是 $s_1$ 和~$s_2$ 的 $x$-变体，且由
      假设有 $s_1(x_i) = s_2(x_i)$。$s_1'(x) = s_2'(x) = m$，这由
      $s_1'$ 和~$s_2'$ 的定义可得。于是归纳
      假设可应用于 $!B$ 和 $s_1'$、$s_2'$，因此
      $\Sat{M}{!B}[s_2']$。故由 $s_2' = \Subst{s_2}{m}{x}$，
      存在某个 $m \in \Domain{M}$ 使得
      $\Sat{M}{!B}[\Subst{s_2}{m}{x}]$，从而
      $\Sat{M}{\indfrm}[s_2]$。}}}

\tagitem{defAll}{}{%
  \iftag{probAll}{%
    \indcase!{!A}{\lforall[x][!B]}{}}{%
    \indcase{!A}{\lforall[x][!B]}{若 $\Sat{M}{\indfrm}[s_1]$，则对
      每个 $m \in \Domain{M}$，都有 $\Sat{M}{!B}[\Subst{s_1}{m}{x}]$。我们
      要证明的是，同样对每个 $m \in \Domain{M}$，
      也有 $\Sat{M}{!B}[\Subst{s_2}{m}{x}]$。为此任取 $m \in \Domain{M}$，
      并考虑 $s_1' = \Subst{s}{m}{x}$ 和 $s_2' =
      \Subst{s}{m}{x}$。我们有 $\Sat{M}{!B}[s_1']$。~$!B$
      的自由变元在 $x_1$, \dots, $x_n$ 和 $x$ 之中。
      $s_1'(x_i) = s_2'(x_i)$，因为 $s_1'$ 和 $s_2'$ 分别是
      $s_1$ 和~$s_2$ 的 $x$-变体，且由假设有
      $s_1(x_i) = s_2(x_i)$。$s_1'(x) = s_2'(x) = m$，这由 $s_1'$ 和~$s_2'$ 的定义可得。于是归纳假设
      可应用于~$!B$ 和~$s_1'$、$s_2'$，因此我们有
      $\Sat{M}{!B}[s_2']$。这对每个 $m \in \Domain{M}$ 都成立，
      即对每个~$m \in \Domain{M}$ 都有 $\Sat{M}{!B}[\Subst{s_2}{m}{x}]$，所以
      $\Sat{M}{\indfrm}[s_2]$。}}}
\end{enumerate}
由归纳法可得，当 $!A$ 的自由变元在 $x_1$, \dots, $x_n$ 之中，
且 $s_1(x_i)=s_2(x_i)$ 对 $i = 1$, \dots,~$n$ 成立时，
有 $\Sat{M}{!A}[s_1]$ 当且仅当 $\Sat{M}{!A}[s_2]$。
\end{proof}

\begin{probtag}{probNot,probOr,probAnd,probIf,probIff,probEx,probAll}
请完成 \olref[fol][syn][ass]{prop:satindep} 的证明。
\end{probtag}

\begin{explain}
语句没有自由变元，因此任意两个变元指派对任何语句的所有（零个）自由变元都赋予相同的值。刚才证明的命题意味着，一个语句在一个结构中是否被满足，完全独立于变元指派的选择。我们要记下这个事实。这为下文给出的（无需提及变元指派的）语句在结构中的满足定义提供了依据。
\end{explain}

\begin{cor}
\ollabel{cor:sat-sentence}
若 $!A$ 是语句，$s$ 是变元指派，则 $\Sat{M}{!A}[s]$ 当且仅当对所有变元指派~$s'$ 都有 $\Sat{M}{!A}[s']$。
\end{cor}

\begin{proof}
  设 $s'$ 为任意变元指派。由于 $!A$ 是语句，它没有自由变元，因此任意变元指派~$s'$ 平凡地与~$s$ 对~$!A$ 的所有自由变元赋予相同的值。于是 \olref{prop:satindep} 的条件得到满足，从而 $\Sat{M}{!A}[s]$ 当且仅当 $\Sat{M}{!A}[s']$。
\end{proof}

\begin{defn}
\ollabel{defn:satisfaction}
若 $!A$ 是语句，我们说结构~$\Struct M$ \emph{满足}~$!A$，记作 $\Sat{M}{!A}$，当且仅当对所有变元指派~$s$ 都有 $\Sat{M}{!A}[s]$。
\end{defn}

若 $\Sat{M}{!A}$，我们也称\emph{$!A$ 在~$\Struct{M}$ 中为真}。满足的概念自然地从单个语句推广到语句的集合。

\begin{defn}
\ollabel{defn:sat}
若 $\Gamma$ 是语句集合，我们说结构~$\Struct M$ \emph{满足}~$\Gamma$，记作 $\Sat{M}{\Gamma}$，当且仅当对所有 $!A \in \Gamma$ 都有 $\Sat{M}{!A}$。
\end{defn}

\begin{prop}\ollabel{prop:sentence-sat-true}
  设 $\Struct{M}$ 为结构，$!A$ 为语句，$s$ 为变元指派。$\Sat{M}{!A}$ 当且仅当 $\Sat{M}{!A}[s]$。
\end{prop}

\begin{proof}
习题。
\end{proof}

\begin{prob}
证明 \olref[fol][syn][ass]{prop:sentence-sat-true}。
\end{prob}

\begin{prop}\ollabel{prop:sat-quant}
  假设 $!A(x)$ 仅含 $x$ 为自由变元，$\Struct M$ 为结构。则：
  \begin{tagenumerate}{prvEx,prvAll}
    \tagitem{prvEx}{$\Sat{M}{\lexists[x][!A(x)]}$ 当且仅当至少存在一个变元指派~$s$ 使得 $\Sat{M}{!A(x)}[s]$。}{}
    \tagitem{prvAll}{$\Sat{M}{\lforall[x][!A(x)]}$ 当且仅当对所有变元指派~$s$ 都有 $\Sat{M}{!A(x)}[s]$。}{}
\end{tagenumerate}
\end{prop}

\begin{proof}
习题。
\end{proof}

\begin{prob}
证明 \olref[fol][syn][ass]{prop:sat-quant}。
\end{prob}

\begin{prob}
\DeclareRobustCommand{\VDash}{\mathrel{||}\joinrel\Relbar}
假设 $\Lang L$ 是没有函数符号的语言。给定结构~$\Struct M$、常项符号 $c$ 以及 $a \in \Domain M$，定义 $\Struct M[a/c]$ 为这样一个结构：除 $\Assign{c}{M[a/c]} = a$ 外，其他与~$\Struct M$ 完全相同。对语句~$!A$，递归定义 $\Struct M \VDash !A$ 如下：
\begin{enumerate}
\tagitem{prvFalse}{%
  \indcase{!A}{\lfalse}{并非 $\Struct{M} \VDash \indfrm$。}}{}

\tagitem{prvTrue}{%
  \indcase{!A}{\ltrue}{$\Struct{M} \VDash \indfrm$。}}{}

\item \indcase{!A}{\Atom{R}{d_1, \dots, d_n}}{$\Struct M \VDash \indfrm$
  当且仅当 $\langle \Assign{d_1}{M}, \dots, \Assign{d_n}{M} \rangle \in
  \Assign{R}{M}$。}

\item \indcase{!A}{\eq[d_1][d_2]}{$\Struct M \VDash \indfrm$ 当且仅当
  $\Assign{d_1}{M} = \Assign{d_2}{M}$。}

\tagitem{prvNot}{%
  \indcase{!A}{\lnot !B}{$\Struct{M} \VDash \indfrm$ 当且仅当
    并非 $\Struct{M} \VDash {!B}$。}}{}

\tagitem{prvAnd}{%
  \indcase{!A}{(!B \land !C)}{$\Struct{M} \VDash
    \indfrm$ 当且仅当 $\Struct{M} \VDash!B$ 且 $\Struct{M} \VDash !C$。}}{}

\tagitem{prvOr}{%
  \indcase{!A}{(!B \lor !C)}{$\Struct{M} \VDash \indfrm$ 当且仅当
    $\Struct{M} \VDash !B$ 或 $\Struct{M} \VDash !C$（或两者皆成立）。}}{}

\tagitem{prvIf}{%
  \indcase{!A}{(!B \lif !C)}{$\Struct{M} \VDash \indfrm$ 当且仅当
    并非 $\Struct {M} \VDash !B$ 或 $\Struct M \VDash !C$（或两者皆成立）。}}{}

\tagitem{prvIff}{%
  \indcase{!A}{(!B \liff !C)}{$\Struct{M} \VDash {\indfrm}$ 当且仅当
    要么 $\Struct{M} \VDash {!B}$ 与 $\Struct{M} \VDash {!C}$ 均成立，要么均不成立。}}{}

\tagitem{prvAll}{%
  \indcase{!A}{\lforall[x][!B]}{$\Struct{M} \VDash {\indfrm}$ 当且仅当
    对所有 $a \in \Domain{M}$，$\Struct{M[a/c]} \VDash \Subst{!B}{c}{x}$，
    前提是 $c$ 不在~$!B$ 中出现。}}{}

\tagitem{prvEx}{%
  \indcase{!A}{\lexists[x][!B]}{$\Struct{M} \VDash
  {\indfrm}$ 当且仅当存在 $a \in \Domain M$ 使得
  $\Struct{M[a/c]} \VDash \Subst{!B}{c}{x}$，前提是 $c$ 不在~$!B$ 中出现。}}{}
\end{enumerate}
设 $x_1$, \dots, $x_n$ 是~$!A$ 中的所有自由变元，
$c_1$, \dots, $c_n$ 是不在~$!A$ 中出现的常项符号，
$a_1$, \dots, $a_n \in \Domain M$，且 $s(x_i) = a_i$。

证明 $\Sat{M}{!A}[s]$ 当且仅当 $\Struct M[a_1/c_1,\dots,a_n/c_n]
\VDash \Subst{\Subst{!A}{c_1}{x_1}\dots}{c_n}{x_n}$。

（这个问题表明，为一阶逻辑给出一种无需变元指派的语义是可能的。）
\end{prob}

\begin{prob}
假设 $f$ 是不在~$!A(x,y)$ 中出现的函数符号。证明：存在结构~$\Struct{M}$ 使得 $\Sat{M}{\lforall[x][\lexists[y][!A(x,y)]]}$ 当且仅当存在结构~$\Struct M'$ 使得 $\Sat{M'}{\lforall[x][!A(x,f(x))]}$。

（本题是 Skolem定理的一个特例；$\lforall[x][!A(x,f(x))]$ 称为 $\lforall[x][\lexists[y][!A(x,y)]]$ 的\emph{Skolem 范式}。）
\end{prob}

\end{document}